% r7le46.tex -- R(7) <= 46 for tournaments (draft)
% Build: tectonic paper/r7le46.tex
% Every number in this file is taken from the result files of the repository
% (results/, data/, notes/); see Section 8 and the comments next to each table.
\documentclass[11pt]{amsart}

\usepackage[margin=1in]{geometry}
\usepackage{amsmath,amssymb,amsthm}
\usepackage{booktabs}
\usepackage{array}
\usepackage{enumitem}
\usepackage{url}
\usepackage[hidelinks]{hyperref}

\setlist{itemsep=2pt,topsep=4pt}
\setcounter{MaxMatrixCols}{10}

\newtheorem{theorem}{Theorem}[section]
\newtheorem{lemma}[theorem]{Lemma}
\newtheorem{proposition}[theorem]{Proposition}
\newtheorem{corollary}[theorem]{Corollary}
\theoremstyle{definition}
\newtheorem{definition}[theorem]{Definition}
\newtheorem{fact}[theorem]{Fact}
\theoremstyle{remark}
\newtheorem{remark}[theorem]{Remark}

\newcommand{\TT}[1]{\mathrm{TT}_{#1}}
\newcommand{\ST}[1]{\mathrm{ST}_{#1}}
\newcommand{\RT}{\mathrm{RT}_5}
\newcommand{\Ostar}{O_{21}^{*}}
\newcommand{\cl}[1]{\textsf{#1}}
\newcommand{\Piv}[2]{\Pi(#1,#2)}
\newcommand{\dom}{\Rightarrow}
\newcommand{\Np}{N^{+}}
\newcommand{\Nm}{N^{-}}
\newcommand{\dpl}{d^{+}}
\newcommand{\dmi}{d^{-}}

\begin{document}

\title{The directed Ramsey number $R(7)$ is at most $46$}

\author{Ahmet Tahir Mat}
\address{Papex Labs}
% TODO: full postal address; no e-mail address in this draft.

\date{Draft, 5 November 2026}

% TODO: check the MSC codes.
\subjclass[2020]{05C20, 05C55, 68V05}
\keywords{Tournament, transitive subtournament, directed Ramsey number,
  SAT solving, cube-and-conquer, LRAT proof checking}

\begin{abstract}
Let $R(k)$ be the least $n$ such that every tournament on $n$ vertices
contains a transitive subtournament on $k$ vertices. In 2022 Neiman, Mackey
and Heule proved $34 \le R(7) \le 47$. We prove $R(7) \le 46$. The argument
looks at the in-degree of a single vertex of a hypothetical counterexample.
It uses three computations. (i)~Complete lists of the tournaments on $21$ to
$25$ vertices that contain no transitive subtournament on $6$ vertices. We
obtain them by a SAT-based enumeration in which every job carries a checked
LRAT certificate; the lists agree with previously published, uncertified
files. (ii)~Five new SAT instances on $46$ vertices with $504$ free arcs each.
We solve them by cube-and-conquer; each of the $48{,}995$ leaf cubes is closed
by an LRAT proof that is checked while it is produced, and the cube cover is
checked as well. (iii)~The $28$ SAT computations of Neiman, Mackey and Heule.
They were published without proofs; we re-ran all of them with the same
checked method ($31{,}942$ leaf cubes). We list explicitly the parts of the
argument that the proof checker does not cover.
\end{abstract}

\maketitle

% =====================================================================
\section{Introduction}\label{sec:intro}

A \emph{tournament} is an orientation of a complete graph. It is
\emph{transitive} if its vertices can be ordered so that every arc points
forward. We write $\TT{k}$ for the transitive tournament on $k$ vertices. Let
$R(k)$ be the least $n$ such that every tournament on $n$ vertices contains
$\TT{k}$ as a subtournament. These numbers go back to Erd\H{o}s and
Moser~\cite{EM64}. The values $R(2),\dots,R(6)$ are $2,4,8,14,28$; see
\cite[\S2]{NMH22} and the references there, in particular~\cite{RP70} for
$R(5)=14$ and~\cite{SF94} for $R(6)=28$.
% TODO: confirm that [RP70] is the right source for R(5) = 14 (the erdosproblems.com
% entry #1216 credits [RePa70] with f(n) >= floor(log2 n + 4 - log2 7) for n >= 14).

For $k=7$ the bounds $32 \le R(7) \le 54$ are due to
Sanchez-Flores~\cite{SF98}.
% TODO: confirm that both 32 and 54 are in [SF98] (NMH cite [5, 9] jointly for 32 <= R(7) <= 53).
Lidick\'y and Pfender lowered the upper bound to $53$ with semidefinite
programming~\cite{LP21}. Neiman, Mackey and Heule~\cite{NMH22} (``NMH'' below)
proved $34 \le R(7) \le 47$ with the help of a SAT solver. Lohn, Lambert and
Heule~\cite{LLH22} later found a further $\TT{7}$-free tournament on
$33$ vertices; the bounds did not change.

NMH's upper bound argument looks at one vertex $v$ of a hypothetical
$\TT{7}$-free tournament $T$. Both neighbourhoods of $v$ are $\TT{6}$-free.
NMH classified the $\TT{6}$-free tournaments on $23$, $24$ and $25$ vertices.
They then used the SAT solver CaDiCaL~\cite{Bie18} to show that certain
partial tournaments have no $\TT{7}$-free completion. Each such partial
tournament consists of a large in-neighbourhood, the vertex $v$, and a small
tournament that must occur in the out-neighbourhood. For $n=47$ this rules out
every in-degree except $22$ and $24$, and a parity argument rules out those
two. For $n=46$ the in/out-degree pairs $(24,21)$ and $(23,22)$ and their
reverses survive, and parity says nothing. NMH remark that a catalogue of the
$\TT{6}$-free tournaments with more than about $20$ vertices ``could likely be
used to further reduce the upper bound on $R(7)$''~\cite[\S6]{NMH22}. We
follow that suggestion.

\begin{theorem}\label{thm:main}
Every tournament on $46$ vertices contains $\TT{7}$. Hence $34 \le R(7) \le 46$.
\end{theorem}

\subsection*{What the proof uses}
The proof (Section~\ref{sec:reduction}) excludes every possible in-degree of
one vertex. Two catalogue facts drive it. Every $\TT{6}$-free tournament on
$22$ vertices contains a fixed $8$-vertex tournament $Y_8$ that NMH used.
Every $\TT{6}$-free tournament on $21$ vertices contains $Y_8$, except the
copies of one particular tournament, which we call $\Ostar$. With these facts, NMH's $25$ instances close the
$(23,22)$ case, and the $(24,21)$ case reduces to five new SAT instances
$\Piv{\cl{24X}}{\Ostar}$, one for each of the five $\TT{6}$-free tournaments on
$24$ vertices. Each new instance has all $46$ vertices and $504$ free arcs.

We say that an unsatisfiability result is \emph{certified} if the SAT solver
emitted a proof in the LRAT format~\cite{CHHKS17} and the checker
\texttt{lrat-check} from the drat-trim distribution~\cite{WHH14} accepted it.
Every unsatisfiability result used in the proof is certified in this sense:
\begin{itemize}
\item the enumeration of the $\TT{6}$-free tournaments on $21$ to $25$
  vertices (Section~\ref{sec:catalogue}), which also re-derives NMH's
  classifications on $23$, $24$ and $25$ vertices;
\item the five new instances (Section~\ref{sec:stage3});
\item all $28$ computations of NMH that the proof needs
  (Section~\ref{sec:stage4}). NMH published their inputs, encoder and solver
  logs but no proofs.
\end{itemize}
The only published theorem that the proof uses without re-deriving it is the
uniqueness of the $\TT{6}$-free tournament on $26$ vertices
(Fact~\ref{fact:st26}). Section~\ref{sec:trust} lists everything else that
lies outside the checked part: the mathematical argument, the encoders, the
checker, and our bookkeeping code.

\subsection*{Prior work}
A literature and code search carried out on 3 November 2026 covered arXiv,
Semantic Scholar and OpenAlex citation data, GitHub, Zenodo, the OEIS and
several lists of open problems. Google Scholar, zbMATH and MathSciNet were
not searched. The search found no improvement of $34 \le R(7) \le 47$ and no
earlier statement of the reduction in Section~\ref{sec:reduction}. This is
a search result, not a claim that nobody has made the observation. The
repository of Lohn, Lambert and Heule contains files of $\TT{6}$-free
tournaments on $20$ to $27$ vertices~\cite{LLHrepo}; no method or
completeness claim is stated for them. Our lists agree with these files for
$21$ to $25$ vertices (Section~\ref{sec:catalogue}).

\subsection*{Organisation}
Section~\ref{sec:prelim} fixes notation. Section~\ref{sec:reduction} proves
Theorem~\ref{thm:main} from a list of named facts. Section~\ref{sec:catalogue}
describes the enumeration of the catalogues and its certification.
Section~\ref{sec:sat} describes the SAT encodings, the cube-and-conquer runs
and the proof checking. Section~\ref{sec:results} reports the results.
Section~\ref{sec:trust} lists the trust base.

% =====================================================================
\section{Preliminaries}\label{sec:prelim}

\subsection{Tournaments}
All tournaments are finite. For distinct vertices $u,v$ of a tournament $T$
exactly one of the arcs $u\to v$, $v\to u$ is present; we then say that $u$
\emph{beats} $v$. For $X \subseteq V(T)$, $T[X]$ is the induced
subtournament. $T$ \emph{contains} $S$ if $T[X] \cong S$ for some $X$.
$T$ is \emph{$\TT{k}$-free} if it does not contain $\TT{k}$. We write
$\Np(v)$ and $\Nm(v)$ for the out- and in-neighbourhood of $v$, identified
with the subtournaments they induce, and $\dpl(v)=|\Np(v)|$,
$\dmi(v)=|\Nm(v)|$. So $\dpl(v)+\dmi(v)=|V(T)|-1$. The \emph{reversal}
$T^{R}$ has every arc reversed. A tournament is transitive if and only if it
has no $3$-cycle.

For tournaments $A$ and $B$ on disjoint vertex sets, $A \dom B$ is the
tournament on $V(A)\cup V(B)$ that contains $A$ and $B$ and in which every
vertex of $A$ beats every vertex of $B$. Let $C_3$ be the $3$-cycle and let
$\RT$ be the regular tournament on $\mathbb{Z}_5$ in which $i \to i+1$ and
$i \to i+2$. Following NMH we use
\[
H_5 = \TT{2} \dom C_3,\quad G_6 = \TT{1} \dom \RT,\quad
H_6 = \TT{3} \dom C_3,\quad G_7 = \TT{2} \dom \RT .
\]
NMH's files call $G_7$ ``W7''. The tournament $Y_8$ is the $8$-vertex
tournament of \cite[Fig.~5]{NMH22}; its adjacency matrix is in
Appendix~\ref{app:data}. It has $15$ three-cycles and is not isomorphic to its
reversal.

$\ST{27}$ is the Paley tournament on $\mathrm{GF}(27)$: $a \to b$ if and only if
$b-a$ is a nonzero square. It is $\TT{6}$-free. All its subtournaments on
$26$ vertices are isomorphic, and so are all its subtournaments on
$25$ vertices (we checked both by computing canonical forms); we call them
$\ST{26}$ and $\ST{25}$. We use NMH's names for the $\TT{6}$-free tournaments on
$24$ and $23$ vertices: $\cl{24A},\dots,\cl{24E}$ are the five isomorphism
classes of $24$-vertex subtournaments of $\ST{27}$, and
$\cl{23A},\dots,\cl{23Y}$ are the $25$ classes on $23$ vertices
($22$ subtournaments of $\ST{27}$ and three doubly regular tournaments).
The tournament $\Ostar$ on $21$ vertices is defined by its digraph6 string in
Appendix~\ref{app:data}. It is regular of out-degree $10$, has $385$
three-cycles and largest transitive subtournament $\TT{5}$, is isomorphic to its
reversal, and has an automorphism group of order $63$. It contains $H_6$ but
neither $G_7$, $Y_8$ nor the reversal of $Y_8$. (None of these properties is
used in the proof except that $\Ostar$ is a fixed tournament.)

\subsection{Pivot instances}
Let $X$ and $Z$ be tournaments. The \emph{pivot instance} $\Piv{X}{Z}$ is the
partial tournament on $V(X) \cup \{p\} \cup V(Z)$ in which the arcs inside $X$
and inside $Z$ are those of $X$ and $Z$, $x \to p$ for $x\in X$, $p \to z$ for
$z \in Z$, and the $|X|\,|Z|$ pairs $\{x,z\}$ with $x \in X$, $z\in Z$ are
\emph{free}. A \emph{completion} orients every free pair. Files and tables name
$\Piv{X}{Z}$ as \texttt{X-p-Z}. Whether $\Piv{X}{Z}$ has a $\TT{7}$-free
completion depends only on the isomorphism classes of $X$ and $Z$.

\begin{lemma}\label{lem:pivot}
Let $T$ be $\TT{7}$-free and $v \in V(T)$. Then $\Nm(v)$ and $\Np(v)$ are
$\TT{6}$-free, and $T^{R}$ is $\TT{7}$-free. If $X' \subseteq \Nm(v)$ and
$Z' \subseteq \Np(v)$, then $T[X' \cup \{v\} \cup Z']$ is a $\TT{7}$-free
completion of a copy of $\Piv{T[X']}{T[Z']}$.
\end{lemma}

\begin{proof}
A $\TT{6}$ in $\Nm(v)$ together with $v$ as its last vertex would be a $\TT{7}$;
similarly, $v$ placed before a $\TT{6}$ in $\Np(v)$ gives a $\TT{7}$. Reversal maps transitive
tournaments to transitive tournaments. The last claim holds because every
vertex of $X'$ beats $v$ and $v$ beats every vertex of $Z'$, and a
subtournament of a $\TT{7}$-free tournament is $\TT{7}$-free.
\end{proof}

\begin{lemma}\label{lem:double}
$R(k) \le 2R(k-1)$ for $k \ge 2$.
\end{lemma}

\begin{proof}
A tournament on $2R(k-1)$ vertices has a vertex $w$ with
$\dpl(w) \ge \lceil (2R(k-1)-1)/2 \rceil = R(k-1)$. So $\Np(w)$ contains
$\TT{k-1}$, and $w$ followed by it is a $\TT{k}$.
\end{proof}

The next lemma and its proof are from \cite[\S5.2.2]{NMH22}.

\begin{lemma}[NMH]\label{lem:h6g7}
Every tournament on at least $20$ vertices contains $\TT{6}$, $H_6$ or $G_7$.
\end{lemma}

\begin{proof}
We begin by showing that every tournament $S$ on at least $10$ vertices contains
$\TT{5}$, $H_5$ or $G_6$. Some vertex $v$ of $S$ has out-degree at least $5$;
let $N$ be a set of $5$ out-neighbours of $v$. If $S[N]$ is regular, then
$S[N] \cong \RT$, so $\{v\}\cup N$ induces $G_6$. (A regular tournament on five
vertices is isomorphic to $\RT$: if $a$ beats $b,c$ with $b \to c$ and $a$ is
beaten by $d,e$ with $d \to e$, then out-degree $2$ forces $c \to d$,
$c \to e$, $b \to d$ and $e \to b$, and $a,b,c,d,e \mapsto 0,1,2,3,4$ is an
isomorphism onto $\RT$.) Otherwise some $u \in N$ beats at least three
vertices of $N$; let $M$ be three of them. If $S[M]$ is transitive, then
$\{v,u\}\cup M$ induces $\TT{5}$. Otherwise $S[M] \cong C_3$ and
$\{v,u\} \cup M$ induces $\TT{2} \dom C_3 = H_5$.

Now let $T$ have at least $20$ vertices. Some vertex $w$ has
$\dpl(w) \ge 10$. Apply the claim to $\Np(w)$. Since $w$ beats all of
$\Np(w)$, $w$ together with a $\TT{5}$, $H_5$ or $G_6$ in $\Np(w)$ induces
$\TT{6}$, $\TT{3}\dom C_3 = H_6$ or $\TT{2} \dom \RT = G_7$.
\end{proof}

% =====================================================================
\section{The reduction for $46$ vertices}\label{sec:reduction}

We state the facts the proof uses, then the lemmas. Facts~\ref{fact:tt5}
and~\ref{fact:cat} are proved by the enumeration of
Section~\ref{sec:catalogue}; Facts~\ref{fact:nmh} and~\ref{fact:new} are the
SAT results of Section~\ref{sec:results}.

\begin{fact}[literature]\label{fact:st26}
Every $\TT{6}$-free tournament on $26$ vertices is isomorphic to $\ST{26}$
\cite{SF94,SF98}; see also \cite[\S2]{NMH22}.
\end{fact}
% TODO: locate the precise statement (paper and theorem number) of the uniqueness of
% ST26 in Sanchez-Flores [SF94]/[SF98]; NMH Sec. 2 states it with citation [8, 9].

\begin{fact}\label{fact:tt5}
There is no $\TT{5}$-free tournament on $14$ vertices; that is,
$R(5) \le 14$. This is classical~\cite{RP70}; our catalogue of
$\TT{5}$-free tournaments (Section~\ref{sec:tt5}) re-derives it.
\end{fact}

\begin{fact}[catalogues]\label{fact:cat}
For $n=21,\dots,25$ let $\mathcal{C}_n$ be the set of isomorphism classes of
$\TT{6}$-free tournaments on $n$ vertices.
\begin{enumerate}[label=(\alph*)]
\item\label{cat25} $\mathcal{C}_{25} = \{\ST{25}\}$.
\item\label{cat24} $\mathcal{C}_{24} = \{\cl{24A},\dots,\cl{24E}\}$.
\item\label{cat23} $\mathcal{C}_{23} = \{\cl{23A},\dots,\cl{23Y}\}$.
\item\label{cat22} Every tournament in $\mathcal{C}_{22}$ contains $Y_8$.
\item\label{cat21} Every tournament in $\mathcal{C}_{21}$ other than $\Ostar$
  contains $Y_8$.
\end{enumerate}
\end{fact}

\begin{fact}[NMH's instances, re-certified here]\label{fact:nmh}
None of the following $28$ pivot instances has a $\TT{7}$-free completion:
$\Piv{\ST{26}}{\TT{5}}$, $\Piv{\ST{25}}{H_6}$, $\Piv{\ST{25}}{G_7}$, and
$\Piv{\cl{23X}}{Y_8}$ for $\cl{X}=\cl{A},\dots,\cl{Y}$.
\end{fact}

\begin{fact}[new instances]\label{fact:new}
For $\cl{X}=\cl{A},\dots,\cl{E}$, the instance $\Piv{\cl{24X}}{\Ostar}$ has no
$\TT{7}$-free completion.
\end{fact}

In the rest of this section $T$ is a $\TT{7}$-free tournament on $46$ vertices
and $v$ is a vertex of $T$. So $\dmi(v)+\dpl(v)=45$.

\begin{lemma}\label{lem:le27}
$\dmi(v) \le 27$.
\end{lemma}

\begin{proof}
By Fact~\ref{fact:tt5} and Lemma~\ref{lem:double}, $R(6) \le 28$. $\Nm(v)$ is
$\TT{6}$-free by Lemma~\ref{lem:pivot}, so it has at most $27$ vertices.
\end{proof}

\begin{lemma}\label{lem:26}
$\dmi(v) \notin \{26,27\}$.
\end{lemma}

\begin{proof}
Suppose $\dmi(v) \in \{26,27\}$. Let $X' \subseteq \Nm(v)$ have $26$ vertices.
$T[X']$ is $\TT{6}$-free, so $T[X'] \cong \ST{26}$ by Fact~\ref{fact:st26}.
$\dpl(v) = 45-\dmi(v) \ge 18$, so by Fact~\ref{fact:tt5} $\Np(v)$ contains a
$\TT{5}$ on some set $Z'$. By Lemma~\ref{lem:pivot},
$\Piv{\ST{26}}{\TT{5}}$ has a $\TT{7}$-free completion. This contradicts
Fact~\ref{fact:nmh}.
\end{proof}

\begin{lemma}\label{lem:25}
$\dmi(v) \ne 25$.
\end{lemma}

\begin{proof}
Suppose $\dmi(v)=25$. Then $\Nm(v) \cong \ST{25}$ by Fact~\ref{fact:cat}\ref{cat25},
and $\Np(v)$ is a $\TT{6}$-free tournament on $20$ vertices. By
Lemma~\ref{lem:h6g7} it contains $H_6$ or $G_7$. By Lemma~\ref{lem:pivot},
$\Piv{\ST{25}}{H_6}$ or $\Piv{\ST{25}}{G_7}$ has a $\TT{7}$-free completion,
contradicting Fact~\ref{fact:nmh}.
\end{proof}

\begin{lemma}\label{lem:23}
$\dmi(v) \ne 23$.
\end{lemma}

\begin{proof}
Suppose $\dmi(v)=23$. By Fact~\ref{fact:cat}\ref{cat23},
$\Nm(v) \cong \cl{23X}$ for some $\cl{X} \in \{\cl{A},\dots,\cl{Y}\}$.
$\Np(v)$ is a $\TT{6}$-free tournament on $22$ vertices, so it contains $Y_8$
by Fact~\ref{fact:cat}\ref{cat22}. Lemma~\ref{lem:pivot} gives a $\TT{7}$-free
completion of $\Piv{\cl{23X}}{Y_8}$, contradicting Fact~\ref{fact:nmh}.
\end{proof}

\begin{lemma}\label{lem:24}
$\dmi(v) \ne 24$.
\end{lemma}

\begin{proof}
Suppose $\dmi(v)=24$. By Fact~\ref{fact:cat}\ref{cat24},
$\Nm(v) \cong \cl{24X}$ for some $\cl{X} \in \{\cl{A},\dots,\cl{E}\}$.
$\Np(v)$ is a $\TT{6}$-free tournament on $21$ vertices. By
Fact~\ref{fact:cat}\ref{cat21} there are two cases.

\emph{$\Np(v)$ contains $Y_8$.} Let $X' \subseteq \Nm(v)$ have $23$ vertices.
$T[X']$ is $\TT{6}$-free, so by Fact~\ref{fact:cat}\ref{cat23} it is isomorphic
to $\cl{23Z}$ for some letter $\cl{Z} \in \{\cl{A},\dots,\cl{Y}\}$.
Lemma~\ref{lem:pivot} gives a $\TT{7}$-free completion of
$\Piv{\cl{23Z}}{Y_8}$, contradicting Fact~\ref{fact:nmh}.

\emph{$\Np(v) \cong \Ostar$.} Then $T = T[\Nm(v) \cup \{v\} \cup \Np(v)]$ is a
$\TT{7}$-free completion of a copy of $\Piv{\cl{24X}}{\Ostar}$, contradicting
Fact~\ref{fact:new}.
\end{proof}

\begin{proof}[Proof of Theorem~\ref{thm:main}]
Suppose $T$ is a $\TT{7}$-free tournament on $46$ vertices and pick any vertex
$v$. Since $\dmi(v)+\dpl(v)=45$, one of the two is at least $23$. If
$\dpl(v) \ge 23$, replace $T$ by $T^{R}$; it is $\TT{7}$-free by
Lemma~\ref{lem:pivot}, and in it $v$ has in-degree $\dpl(v)$. So we may assume
$\dmi(v) \ge 23$. Lemmas~\ref{lem:le27}--\ref{lem:24} exclude every value
$\dmi(v) \ge 23$. The lower bound $R(7) \ge 34$ is from~\cite{NMH22}.
\end{proof}

Table~\ref{tab:deps} lists which facts each step uses.

\begin{table}[ht]
\centering\small
\caption{Facts used by each step of the proof of Theorem~\ref{thm:main}.}
\label{tab:deps}
\begin{tabular}{@{}lll@{}}
\toprule
Step & Case & Facts used \\
\midrule
Lemma~\ref{lem:le27} & $\dmi(v) \ge 28$ & Fact~\ref{fact:tt5}, Lemma~\ref{lem:double} \\
Lemma~\ref{lem:26} & $\dmi(v)\in\{26,27\}$ & Facts~\ref{fact:st26}, \ref{fact:tt5}; $\Piv{\ST{26}}{\TT{5}}$ \\
Lemma~\ref{lem:25} & $\dmi(v)=25$ & Fact~\ref{fact:cat}\ref{cat25}, Lemma~\ref{lem:h6g7}; $\Piv{\ST{25}}{H_6}$, $\Piv{\ST{25}}{G_7}$ \\
Lemma~\ref{lem:23} & $\dmi(v)=23$ & Fact~\ref{fact:cat}\ref{cat23},\ref{cat22}; $25\times\Piv{\cl{23X}}{Y_8}$ \\
Lemma~\ref{lem:24} & $\dmi(v)=24$ & Fact~\ref{fact:cat}\ref{cat24},\ref{cat23},\ref{cat21}; $25\times\Piv{\cl{23X}}{Y_8}$; $5\times\Piv{\cl{24X}}{\Ostar}$ \\
Theorem~\ref{thm:main} & $\dmi(v) \le 22$ & reversal (Lemma~\ref{lem:pivot}) \\
\bottomrule
\end{tabular}
\end{table}

\begin{remark}
The argument needs one vertex and no counting. For $n=46$ counting alone could
not work: in-degree sequences with $23$ vertices of in-degree $21$ and $23$ of
in-degree $24$, or $23$ of each of $22$ and $23$, sum to
$\binom{46}{2}=1035$.
\end{remark}

\begin{remark}\label{rem:h6sat}
The catalogue is needed in the case $\dmi(v)=24$. By Lemma~\ref{lem:h6g7},
$\Np(v)$ contains $H_6$ or $G_7$, but this is not enough: in exploratory runs
(not part of the proof) the solver found $\TT{7}$-free completions of
$\Piv{\cl{24X}}{H_6}$ for $\cl{X}=\cl{A},\cl{B},\cl{C}$. We re-checked each of
these $31$-vertex tournaments independently: it is a completion of the
instance and its largest transitive subtournament has $6$ vertices.
\end{remark}

% =====================================================================
\section{The catalogues and their certification}\label{sec:catalogue}

\subsection{$\TT{5}$-free tournaments}\label{sec:tt5}
We built the $\TT{5}$-free tournaments on $n \le 14$ vertices by one-vertex
extension. Every $\TT{5}$-free tournament on $n+1$ vertices minus a vertex is
$\TT{5}$-free on $n$ vertices. So the list for $n+1$ is obtained from the list
for $n$ by adding a vertex in all $2^n$ ways, keeping the $\TT{5}$-free
results, and reducing them to isomorphism classes with nauty's \texttt{labelg}
\cite{MP14}. The numbers of classes for $n=1,\dots,14$ are
\[
1,\,1,\,2,\,4,\,11,\,35,\,121,\,376,\,692,\,369,\,45,\,1,\,1,\,0 .
\]
For $n \le 10$ the lists equal, class by class, the output of nauty's
\texttt{gentourng} filtered for $\TT{5}$. For $n \le 13$ they equal the
$\TT{5}$-free files in NMH's repository~\cite{NMHrepo}. The value $0$ at
$n=14$ is Fact~\ref{fact:tt5}. This computation is exhaustive but it is not
certified by a proof checker (Section~\ref{sec:trust}).

\subsection{Pivot decomposition}

\begin{proposition}\label{prop:pivotdec}
Let $T$ be a $\TT{6}$-free tournament on $n$ vertices and let $w$ be a vertex of
maximum out-degree $d$. Then $\lceil (n-1)/2 \rceil \le d \le 13$ and
$n-1-d \le 13$. Let $A$ and $B$ be the members of the $\TT{5}$-free lists
isomorphic to $\Nm(w)$ and $\Np(w)$. Then $T$ is isomorphic to a completion
of $\Piv{A}{B}$ in which every vertex has out-degree at most $d$.
\end{proposition}

\begin{proof}
The average out-degree is $(n-1)/2$. $\Nm(w)$ and $\Np(w)$ are $\TT{5}$-free
(Lemma~\ref{lem:pivot} with $6$ in place of $7$), so they have at most $13$
vertices by Fact~\ref{fact:tt5}. Relabel $T$ along isomorphisms
$\Nm(w) \to A$ and $\Np(w) \to B$; the result is a completion of $\Piv{A}{B}$,
and $d$ bounds all out-degrees because it is the maximum.
\end{proof}

For each $n$, a \emph{job} is a triple $(d,i,j)$ with
$\lceil (n-1)/2 \rceil \le d \le 13$, $A_i$ in the $\TT{5}$-free list on
$n-1-d$ vertices and $B_j$ in the list on $d$ vertices. The job asks for all
$\TT{6}$-free completions of $\Piv{A_i}{B_j}$ with all out-degrees at most $d$.
There are $46$, $414$, $3{,}086$, $17{,}673$ and $167{,}798$ jobs for
$n=25,24,23,22,21$. For example, for $n=21$ the values $d=10,11,12,13$ give
$369\cdot 369$, $692\cdot 45$, $376 \cdot 1$ and $121\cdot 1$ jobs.

\subsection{The formula of a job}
The variables of the job $(d,i,j)$ are one Boolean variable per free pair
$\{x,y\}$, $x \in A_i$, $y\in B_j$, true if and only if $x \to y$, plus
auxiliary variables. The clauses are:
\begin{enumerate}[label=(\roman*)]
\item the $\TT{6}$-exclusion clauses of Section~\ref{sec:encoding}. With
  $k=6$ every $6$-set has at most $6!=720$ orderings, below the limit used
  ($5900$), so all clauses are of the ``direct'' kind;
\item cardinality constraints in Sinz's sequential-counter
  encoding~\cite{Sinz05}: each $x \in A_i$ beats at most
  $d-1-\dpl_{A_i}(x)$ vertices of $B_j$, and each $y \in B_j$ beats at most
  $d-\dpl_{B_j}(y)$ vertices of $A_i$.
\end{enumerate}
Every $\TT{6}$-free completion with all out-degrees at most $d$ satisfies the
formula, with the counter variables set to the partial sums
(Lemma~\ref{lem:sound} and~\cite{Sinz05}).

\subsection{Enumeration and certification}
The enumerator (\path{src/pivenum.cpp}) solves each job with the CaDiCaL
library in incremental mode. After each model it checks that the completed
tournament is $\TT{6}$-free and respects the degree bound, records it, and adds
a \emph{blocking clause} that excludes its assignment of the arc variables.

A job is \emph{certified} when the job formula plus one blocking clause per
recorded model is unsatisfiable and this is certified by a checked LRAT proof
(\path{scripts/certify_enum.py}). Certification also re-checks, in separate
code, that every recorded model is $\TT{6}$-free. A certified job therefore
has no completion, in the sense of the job, other than the recorded ones.

For $n=25$, $24$ and $22$ every job was enumerated and then certified. For
$n=23$ and $21$ we used a shortcut (\path{scripts/certify_first.py}). In a
sample of $168$ jobs for $n=21$ it needed $42\%$ of the solver time of the
two-step path. Each job formula is solved once with
LRAT output and the proof is checked. If the answer is unsatisfiable and the
proof is accepted, the job has no models, and this one checked solve is both
its enumeration and its certificate. Otherwise the job is enumerated as above
and later certified with its blocking clauses. Of the $167{,}798$ jobs for
$n=21$, $167{,}632$ were closed by the single checked solve and $166$ by the
blocking-clause certificate; for $n=23$ the numbers are $3{,}052$ and $34$.

The run for $n=21$ was interrupted once when its cloud instance was reclaimed.
It was resumed on another instance from the synchronised records. In the
synchronised files one job ($d=11$) had a completion record but no certificate
record. The certification step treats every job without a certificate record
as uncertified, and it certified this job again (unsatisfiable, proof
accepted, no models).

\begin{proposition}\label{prop:complete}
Assume that the $\TT{5}$-free lists are complete, that every
$\TT{6}$-free completion of a job with out-degrees at most $d$ satisfies the
job formula, and that each recorded model agrees with the fixed arcs of its
job. If every job for $n$ is certified, then every $\TT{6}$-free tournament on
$n$ vertices is isomorphic to a recorded model.
\end{proposition}

\begin{proof}
Let $T$ be $\TT{6}$-free on $n$ vertices. By
Proposition~\ref{prop:pivotdec}, $T$ is isomorphic to a completion $T'$ of
$\Piv{A_i}{B_j}$ with out-degrees at most $d$, for some job $(d,i,j)$. The arc
assignment of $T'$ extends to a model of the job formula. Since the job
formula plus the blocking clauses is unsatisfiable, the arc assignment of $T'$
is blocked, so it equals the arc assignment of a recorded model $M$. $T'$ and
$M$ agree on the fixed arcs and on the free arcs, so $T' = M$.
\end{proof}

\subsection{Classes and checks}
The recorded models are reduced to isomorphism classes with \texttt{labelg}.
Table~\ref{tab:cat} gives the numbers. The following checks were made on the
class representatives (\path{scripts/stage12_check.py}).
\begin{itemize}
\item For each $n = 21,\dots,25$, the set of canonical forms equals the set of
  canonical forms of the file \texttt{tt6free$n$.d6} of~\cite{LLHrepo}.
\item For $n=23,24,25$ the classes equal NMH's: our $\ST{25}$ (from the Paley
  construction), our $\cl{24A},\dots,\cl{24E}$ and $22$ sub-$\ST{27}$ classes
  on $23$ vertices (matched to NMH's letters by canonical form), the doubly
  regular $\cl{23W}$ and $\cl{23X}$ from NMH's repository (originally from
  McKay's data~\cite{McKayDigraphs}), and $\cl{23Y}$, which is the
  quadratic-residue tournament on $\mathbb{Z}_{23}$.
\item All $89$ classes on $22$ vertices contain $Y_8$.
\item All $17{,}040$ classes on $21$ vertices except one contain $Y_8$, and the
  exception has the canonical form of $\Ostar$.
\end{itemize}
Containment of $Y_8$ is decided by a backtracking search for an induced
embedding (\path{src/tourn.py}). Together with
Proposition~\ref{prop:complete} these checks give Fact~\ref{fact:cat}. Only one
property of \texttt{labelg} is needed for this: its output is a relabelling
of its input. Equal canonical forms then imply isomorphism, and every recorded
model is isomorphic to its class representative. The class counts and the
comparison with the other catalogues also rely on canonicity.

\begin{table}[ht]
\centering\small
\caption{Enumeration of the $\TT{6}$-free tournaments on $n$ vertices.
``Jobs with models'' counts jobs with at least one $\TT{6}$-free completion;
``models'' counts labelled completions. Solver hours are sums of per-job
solver wall-clock times, enumeration and certification together. ``Single
checked solve'': jobs without models were closed by one checked solve; the
others were enumerated and certified with blocking clauses.}
% Source: results/enum/k6n{21..25}/{summary,certify_summary}.json, records.jsonl,
% certify.jsonl, certify_first.jsonl; results/stage12_check.json.
\label{tab:cat}
\begin{tabular}{@{}rrrrrlr@{}}
\toprule
$n$ & jobs & jobs with models & models & classes & certification path & solver h \\
\midrule
25 & 46 & 2 & 78 & 1 & enumerate, then certify & 0.1 \\
24 & 414 & 10 & 390 & 5 & enumerate, then certify & 1.4 \\
23 & 3{,}086 & 34 & 1{,}068 & 25 & single checked solve & 5.0 \\
22 & 17{,}673 & 62 & 2{,}098 & 89 & enumerate, then certify & 507.5 \\
21 & 167{,}798 & 165 & 200{,}677 & 17{,}040 & single checked solve & 1{,}146.4 \\
\bottomrule
\end{tabular}
\end{table}

% =====================================================================
\section{SAT encodings, cube-and-conquer and proof checking}\label{sec:sat}

\subsection{The $\TT{k}$-exclusion encoding}\label{sec:encoding}
Let $P$ be a partial tournament on vertices $0,\dots,n-1$: each pair is a
fixed arc or free. Our encoder (\path{src/encode.c}; \path{src/ttenc.h} implements the same
construction for the enumerator) produces the following CNF $F_L(P)$
for $\TT{k}$-exclusion, with a parameter $L$.
\begin{itemize}
\item \emph{Arc variables.} One variable $x_{ij}$ per free pair $i<j$, true if
  and only if $i \to j$.
\item \emph{$k$-sets.} For every $k$-set $S$ that contains no $3$-cycle of fixed
  arcs, list the orderings of $S$ compatible with the fixed arcs, that is, the
  transitive orientations of $S$ that agree with them. If there are none, $S$
  is skipped. If there are at most $L$, the set of assignments of the free
  pairs inside $S$ that make $S$ transitive is reduced to its prime
  implicants by Quine--McCluskey merging, and one clause excluding each prime
  implicant is emitted (``direct'' clauses). If there are more than $L$, one
  ``$3$-cycle clause'' is emitted: the disjunction, over the triples of $S$,
  of a literal saying that the triple is cyclic.
\item \emph{$3$-cycle literals.} For a triple with no free pair the literal is
  constant. For a triple with exactly one free pair whose fixed arcs form a
  path $a \to b \to c$, it is the arc literal for $c \to a$; if the two fixed
  arcs do not form a path, the triple cannot be cyclic and the literal is
  false. For a triple with
  two or more free pairs it is a fresh variable $t_{abc}$ with the clauses
  $\neg t_{abc} \vee \neg[u\to w] \vee [w\to y]$ for the three rotations
  $(u,w,y)$ of each of the two cyclic orientations of $\{a,b,c\}$; clauses
  that are satisfied by fixed arcs are dropped.
\end{itemize}

\begin{lemma}[soundness]\label{lem:sound}
Let $T$ be a $\TT{k}$-free completion of $P$. Set $x_{ij}$ to the orientation
of $\{i,j\}$ in $T$ and $t_{abc}$ to true if and only if $\{a,b,c\}$ is a
$3$-cycle in $T$. This assignment satisfies $F_L(P)$ for every $L$. Hence, if
$F_L(P)$ is unsatisfiable, $P$ has no $\TT{k}$-free completion.
\end{lemma}

\begin{proof}
A clause $\neg t_{abc} \vee \neg[u\to w] \vee [w\to y]$ can only fail if
$\{a,b,c\}$ is cyclic and $u \to w$; then the cycle is $u \to w \to y \to u$, so
$w \to y$. A direct clause excludes a set of assignments of the free pairs of
$S$ all of which make $S$ transitive, because Quine--McCluskey only merges
cubes contained in that set. $T[S]$ is not transitive, so its assignment is
not excluded. For a $3$-cycle clause, $T[S]$ is not transitive, so it has a
cyclic triple. That triple is not fixed-cyclic (such sets are skipped), so its
literal is an arc literal or a variable $t_{abc}$, and it is true.
\end{proof}

Lemma~\ref{lem:sound} is a statement about the construction. That the program
implements it is a separate matter (Section~\ref{sec:trust}). The encoder was
compared with brute force on $60$ random partial tournaments with $8$ to $10$
vertices, $k\in\{4,5\}$, up to $13$ free pairs and $L \in \{1,3,10,5900\}$.
The CNF was evaluated on all $74{,}944$ assignments of the free pairs, with the
$t$-variables set as in Lemma~\ref{lem:sound}; it agreed with the brute-force
answer every time (\path{tests/test_encode.py}).

\subsection{The five new instances}\label{sec:stage3enc}
Each instance $\Piv{\cl{24X}}{\Ostar}$ was built from the matrices
\path{data/I24_X.txt} and \path{data/O21star.txt} (Appendix~\ref{app:data})
in the vertex order: $\cl{24X}$ (vertices $0$--$23$), $p$ ($24$), $\Ostar$
($25$--$45$). It was
encoded with $k=7$ and $L=1$. Every $7$-set with exactly one compatible
ordering gets one direct clause; all others get a $3$-cycle clause. Each
formula has $504$ arc variables and $10{,}836$ variables $t_{abc}$, one for
every triple with two vertices on one side of the pivot and one on the other.
For $\cl{24A}$ there are $2{,}608{,}200$ direct clauses, $6{,}355{,}503$
$3$-cycle clauses and $43{,}344$ definition clauses, $9{,}007{,}047$ clauses and
$208{,}977{,}552$ literals in all. The other four formulas are of the same size
(Table~\ref{tab:stage3}).

\subsection{NMH's instances and the clause checker}\label{sec:stage4enc}
For the $28$ instances of Fact~\ref{fact:nmh} we did not use our encoder. We
used NMH's encoder \path{encoding/tour-encode.c}~\cite{NMHrepo}, compiled as in
their Makefile, on NMH's own input matrices for the $27$ instances that have a
log, with the limit $5040 = 7!$. With this limit every $7$-set gets direct
clauses and there are no auxiliary variables. For $\Piv{\ST{26}}{\TT{5}}$,
for which NMH published no input, the input is $\ST{27}$ minus a vertex, then
$p$, then $\TT{5}$. For all $27$ logged instances, the header, the number of
clauses and the size in bytes of our CNF equal those reported in NMH's
CaDiCaL logs. NMH's CNF files are not published, so this is evidence that we
solve the formulas NMH solved, not a byte-for-byte comparison. The inputs, the
encoder source and all $28$ CNFs are pinned by SHA-256 in
\path{data/stage4_manifest.json}; the cloud driver rebuilt every CNF and
refused any mismatch.

We do not rely on NMH's encoder for soundness. A separate program
(\path{src/cnf_sound.c}, about $230$ lines of C) checks, for a partial
tournament $P$ and a CNF $F$ over its arc variables, that every clause of $F$
holds in every $\TT{7}$-free completion of $P$. It accepts a clause $C$ if
\begin{enumerate}[label=(\alph*)]
\item some literal of $C$ is true on a fixed arc of $P$ (this covers NMH's
  unit clauses for the fixed arcs, which must agree with $P$); or
\item $C$ contains a literal and its negation; or
\item after the literals that are false on fixed arcs are removed, let
  $\alpha$ be the orientation of the remaining free pairs that falsifies $C$;
  there is a $7$-set $S$ containing every endpoint of $\alpha$ such that every
  orientation of the free pairs of $S$ that extends $\alpha$ makes $S$
  transitive.
\end{enumerate}
It also checks the clause count and the variable range against the header.
In case~(c) every completion that falsifies $C$ contains a $\TT{7}$ on $S$, so
$C$ holds in every $\TT{7}$-free completion. Hence, if the checker accepts all
clauses and $F$ is unsatisfiable, $P$ has no $\TT{7}$-free completion,
whatever program produced $F$. The checker accepted all $28$ CNFs with no
failure. Its tests (\path{tests/test_cnf_sound.py}) compare it with brute force
on $40$ random partial tournaments; of $240$ mutated clauses, $42$ were unsound
by brute force, and the checker accepted none of them.

Finally a structure check (\path{scripts/stage4_instances.py}) confirms that
each input is a pivot instance with exactly the $X \times Z$ pairs free, that
$X$ is isomorphic to our own representative of the class named by NMH (from
our constructions, Section~\ref{sec:catalogue}), that the $25$ classes
$\cl{23A},\dots,\cl{23Y}$ are pairwise non-isomorphic, and that $Z$ is
isomorphic to our $Y_8$, $H_6$, $G_7$ or $\TT{5}$.

\subsection{Cube-and-conquer}\label{sec:cnc}
All $33$ instances were solved by cube-and-conquer~\cite{HKWB11} with the same
scripts (\path{scripts/make_cubes.py}, \path{scripts/run_cubes.py}).
\begin{itemize}
\item \emph{Level $0$.} The $K=10$ arc variables that occur most often in the
  formula are chosen, and the $2^{10}=1024$ cubes are all sign patterns on
  them. In each new instance the ten variables are the arcs between one vertex
  of $\cl{24X}$ and ten vertices of $\Ostar$.
\item \emph{Conquer.} A cube $c$ is solved as $F \wedge c$ ($F$ plus one unit
  clause per literal of $c$) by CaDiCaL with LRAT output and a time limit that
  depends on the level.
\item \emph{Splitting.} A cube that reaches its time limit is replaced by
  $2^{S}=16$ children: the cube extended by all sign patterns on the $S=4$ most
  frequent arc variables that it does not fix yet. The children run at the
  next level. The last level has no time limit.
\item \emph{SAT.} A satisfiable cube would stop the run and its model would be
  checked independently (\path{scripts/verify_model.py}). No cube was
  satisfiable.
\end{itemize}

\subsection{Streamed LRAT checking}\label{sec:lrat}
Text LRAT proofs grew at $2$ to $6$ megabytes per second of solving in our
measurements, so proofs for thousands of solver-hours cannot be stored. Every
cube is therefore checked while it is solved (\path{scripts/stream_solve.py}). The
script creates three named pipes. One carries the formula $F\wedge c$ to
CaDiCaL. One carries a second copy of the same formula, written by the same
code from the same header, clause file and unit clauses, to
\texttt{lrat-check}. The third carries CaDiCaL's text LRAT proof to
\texttt{lrat-check}. The checker thus reads the formula itself; it does not
take it from the solver. A cube is recorded as closed only if CaDiCaL reports
unsatisfiable and the checker prints \texttt{c VERIFIED}. \texttt{lrat-check}
prints this only if it derived the empty clause and found no error. The proof
is then discarded. Only the per-cube record is kept: identifier, literals,
level, status, checker verdict and times.

\subsection{Cover check}
At the end, \path{scripts/cnc_summary.py} checks three things for each
instance.
\begin{enumerate}[label=(\arabic*)]
\item \emph{Level-$0$ cover.} The formula $\bigwedge_{c} \neg c$ over the $1024$
  level-$0$ cubes is unsatisfiable. This is shown by CaDiCaL with a checked LRAT
  proof, from the cubes alone, without the instance formula.
\item \emph{Splits.} For every cube with children, the children are exactly
  the $2^{S}$ sign patterns on one common sequence of $S$ variables not fixed by
  the parent, each appended to the parent's literals.
\item \emph{Leaves.} Every cube without children is recorded as unsatisfiable
  with checker verdict \texttt{VERIFIED}; no cube is open and none is
  satisfiable.
\end{enumerate}
It prints \texttt{ALL CLOSED} only if all three hold.

\begin{lemma}\label{lem:cover}
If (1)--(3) hold, then $F$ is unsatisfiable.
\end{lemma}

\begin{proof}
Let $\beta$ be an assignment. By (1), $\beta$ satisfies some level-$0$ cube.
If a cube satisfied by $\beta$ has children, then by (2) $\beta$ satisfies the
child whose sign pattern agrees with $\beta$ on the split variables. The cubes
get longer, so this descent ends at a leaf $c$ satisfied by $\beta$. By (3),
$F \wedge c$ is unsatisfiable, so $\beta$ falsifies $F$.
\end{proof}

Lemma~\ref{lem:cover} does not depend on the time limits, on the depth of the
splits, or on the order in which cubes were run. This matters for the run history of NMH's instances in
Section~\ref{sec:stage4}.

\subsection{Tools}\label{sec:tools}
We used CaDiCaL~3.0.1~\cite{Bie18,CaDiCaL2}, built from the release tarball
(SHA-256 \texttt{0a8ea563\ldots{}7e11}). The checker \texttt{lrat-check} is the
file \path{lrat-check.c} of the drat-trim distribution, with the header
``Copyright (c) 2017--2023 Marijn Heule'' and ``last edit July 2023''
(SHA-256 \texttt{bf07c2ac\ldots{}4a95}). It was compiled with
\texttt{gcc -O2 -std=c99 -DLONGTYPE}.
% TODO: record the upstream drat-trim commit from which lrat-check.c was taken.
nauty~2.9.3~\cite{MP14} was used for \texttt{labelg} and \texttt{gentourng}
(tarball SHA-256 \texttt{9fc4edae\ldots{}1b5b}). The scripts are in Python~3.
The cloud instances ran Ubuntu 24.04 (arm64). The development machine (Apple
M5) had the same versions of CaDiCaL and nauty.

% =====================================================================
\section{Results}\label{sec:results}

\subsection{Catalogues}
All jobs for $n=21,\dots,25$ were completed and certified. Table~\ref{tab:cat}
gives the counts. The classes equal the files of~\cite{LLHrepo}, and for
$n=23,24,25$ also NMH's classes. Fact~\ref{fact:cat} holds.

\subsection{The five new instances}\label{sec:stage3}
All five instances are \texttt{ALL CLOSED} (Table~\ref{tab:stage3}). The time
limits were $600$\,s at level $0$ and $1800$\,s at level $1$, with no limit at
level $2$. Of the $5 \times 1024$ level-$0$ cubes, $2{,}195$ were closed
directly and $2{,}925$ were split into $16$ children each. All $46{,}800$
level-$1$ cubes were closed within the level-$1$ limit, so level $2$ was never
used. In total $48{,}995$ leaves were closed by checked proofs; no cube was
satisfiable, no check failed, and no leaf is open. This proves
Fact~\ref{fact:new}.

The five instances ran side by side on one cloud instance with $96$ vCPUs, with
$15$ workers each, from 18:42 UTC on 3 November to 06:07 UTC on 5 November 2026.
The summaries were recomputed on a second machine from the synchronised
records, including the level-$0$ cover check, with the same result
(\path{summary.mac.json}). This recomputation does not re-check the leaves,
whose proofs were not kept.

\begin{table}[ht]
\centering\small
\caption{The five new instances $\Piv{\cl{24X}}{\Ostar}$ ($46$ vertices,
$504$ arc variables, $11{,}340$ variables in all). ``Split vertex'' is the
vertex of $\cl{24X}$ (numbered from $0$) whose arcs carry the level-$0$ split.
``L0'', ``L1'': level $0$, level $1$. ``Max s'' is the longest single cube in
seconds.}
% Source: results/stage3/24X-p-O21s/run/{summary.json,results.jsonl}, results/stage3-ABCDE.log.
\label{tab:stage3}
\setlength{\tabcolsep}{4pt}
\begin{tabular}{@{}lrcrrrrrr@{}}
\toprule
Instance & clauses & split vertex & L0 closed & L0 split & L1 cubes & leaves & max s & solver h \\
\midrule
\texttt{24A-p-O21*} & 9{,}007{,}047 & 21 & 434 & 590 & 9{,}440 & 9{,}874 & 1{,}330 & 510.8 \\
\texttt{24B-p-O21*} & 8{,}988{,}168 & 22 & 429 & 595 & 9{,}520 & 9{,}949 & 1{,}402 & 523.7 \\
\texttt{24C-p-O21*} & 8{,}988{,}168 & 3 & 461 & 563 & 9{,}008 & 9{,}469 & 1{,}431 & 513.5 \\
\texttt{24D-p-O21*} & 9{,}007{,}047 & 8 & 439 & 585 & 9{,}360 & 9{,}799 & 1{,}683 & 525.8 \\
\texttt{24E-p-O21*} & 9{,}007{,}047 & 20 & 432 & 592 & 9{,}472 & 9{,}904 & 1{,}544 & 530.2 \\
\midrule
Total & & & 2{,}195 & 2{,}925 & 46{,}800 & 48{,}995 & & 2{,}604.1 \\
\bottomrule
\end{tabular}
\end{table}

\subsection{NMH's instances}\label{sec:stage4}
All $28$ instances are \texttt{ALL CLOSED} (Table~\ref{tab:stage4}), with
$31{,}942$ leaves closed by checked proofs, no satisfiable cube and no failed
check. This proves Fact~\ref{fact:nmh} without relying on NMH's runs. The
summaries recomputed on a second machine are identical for all $28$ instances.

\emph{Run history.} The run used two successive cloud instances with $96$ vCPUs
each. On the original machine, from 12:21 UTC on 4 November, the driver ran six instances
at a time with $15$ workers each and time limits $(1800, 3600, \infty)$ seconds
for levels $0$, $1$, $2$. At 18:00 UTC a second driver process was started on
the same machine for the $22$ instances that had not started yet, with $4$
workers each and the same limits; the two drivers worked on disjoint
instances. At about 01:27 UTC on 5 November the original machine was lost to a
spot interruption. The run was resumed at 02:12 UTC on a second machine from the
synchronised records, with all $28$ instances side by side, $6$ workers each,
and the time limits changed to $(1800, 3600, 1800, 7200, \infty)$ for levels $0$
to $4$. Eight instances were already closed. The last instance closed at 10:03
UTC.
% NOTE (for the author): the commit message of 7bae8ed says "two spot reclaims". The
% repository logs show one spot interruption of a Stage 4 machine (01:27 UTC, 5 Nov)
% and one of the enumeration machine (about 10:16 UTC, 4 Nov); the 18:00 UTC event was
% an additional driver on the same machine. Please confirm this account.

The change of time limits does not affect correctness. Time limits only
decide which cubes are split. A cube recorded as closed stays a leaf. A cube
recorded as timed out has children. Work that was running at the interruption
left no record and was redone. Lemma~\ref{lem:cover} needs only the checks
(1)--(3), which \path{cnc_summary.py} applies to the final records whatever
the depth: the children of every split must be all $2^{S}$ sign patterns on
fresh variables, and every leaf must carry its own accepted proof. The deepest
cubes are at level $4$ ($26$ literals).

\begin{table}[ht]
\centering\small
\caption{NMH's instances, rebuilt with NMH's encoder and re-run with checked
LRAT proofs. ``Free'': free arcs. NMH columns: CaDiCaL version, run time in
hours and conflicts in millions, from NMH's logs (no log was published for
$\Piv{\ST{26}}{\TT{5}}$). ``Splits'': cubes that timed out and were split;
``depth'': deepest level.}
% Source: data/stage4_manifest.json (sizes, NMH log data), results/stage4/NAME/run/{summary.json,results.jsonl}.
\label{tab:stage4}
\setlength{\tabcolsep}{3.5pt}
\begin{tabular}{@{}lrrcrrrrrr@{}}
\toprule
 & & & \multicolumn{3}{c}{NMH (logs)} & \multicolumn{4}{c}{this work} \\
\cmidrule(lr){4-6}\cmidrule(l){7-10}
Instance & free & clauses & CaDiCaL & h & $10^6$ confl. & splits & depth & leaves & solver h \\
\midrule
\texttt{ST26-p-TT5} & 130 & 1{,}099{,}282 & -- & -- & -- & 0 & 0 & 1{,}024 & 5.6 \\
\texttt{ST25-p-H6} & 150 & 1{,}687{,}063 & 1.5.2 & 23.9 & 88 & 1 & 1 & 1{,}039 & 10.8 \\
\texttt{ST25-p-G7} & 175 & 2{,}603{,}275 & 1.5.2 & 33.5 & 96 & 0 & 0 & 1{,}024 & 11.4 \\
\texttt{23A-p-Y8} & 184 & 2{,}353{,}250 & 1.3.1 & 121.0 & 354 & 12 & 4 & 1{,}204 & 53.8 \\
\texttt{23B-p-Y8} & 184 & 2{,}364{,}374 & 1.3.1 & 79.4 & 265 & 12 & 4 & 1{,}204 & 47.9 \\
\texttt{23C-p-Y8} & 184 & 2{,}364{,}374 & 1.3.1 & 97.5 & 308 & 12 & 4 & 1{,}204 & 47.5 \\
\texttt{23D-p-Y8} & 184 & 2{,}353{,}250 & 1.3.1 & 157.1 & 427 & 7 & 4 & 1{,}129 & 50.1 \\
\texttt{23E-p-Y8} & 184 & 2{,}343{,}727 & 1.3.1 & 203.9 & 543 & 8 & 4 & 1{,}144 & 53.0 \\
\texttt{23F-p-Y8} & 184 & 2{,}343{,}727 & 1.3.1 & 230.9 & 609 & 8 & 4 & 1{,}144 & 52.7 \\
\texttt{23G-p-Y8} & 184 & 2{,}343{,}727 & 1.3.1 & 212.2 & 560 & 11 & 4 & 1{,}189 & 55.4 \\
\texttt{23H-p-Y8} & 184 & 2{,}343{,}727 & 1.3.1 & 249.9 & 650 & 10 & 4 & 1{,}174 & 55.0 \\
\texttt{23I-p-Y8} & 184 & 2{,}353{,}250 & 1.3.1 & 203.7 & 495 & 7 & 4 & 1{,}129 & 46.0 \\
\texttt{23J-p-Y8} & 184 & 2{,}353{,}250 & 1.3.1 & 240.0 & 589 & 5 & 2 & 1{,}099 & 44.9 \\
\texttt{23K-p-Y8} & 184 & 2{,}353{,}250 & 1.3.1 & 138.5 & 383 & 9 & 4 & 1{,}159 & 48.5 \\
\texttt{23L-p-Y8} & 184 & 2{,}343{,}727 & 1.3.1 & 242.6 & 618 & 9 & 4 & 1{,}159 & 53.8 \\
\texttt{23M-p-Y8} & 184 & 2{,}364{,}374 & 1.3.1 & 94.3 & 300 & 10 & 4 & 1{,}174 & 46.5 \\
\texttt{23N-p-Y8} & 184 & 2{,}364{,}374 & 1.2.1 & 235.5 & 390 & 10 & 4 & 1{,}174 & 45.9 \\
\texttt{23O-p-Y8} & 184 & 2{,}353{,}250 & 1.2.1 & 310.1 & 460 & 12 & 4 & 1{,}204 & 52.3 \\
\texttt{23P-p-Y8} & 184 & 2{,}364{,}374 & 1.2.1 & 198.8 & 342 & 10 & 4 & 1{,}174 & 45.4 \\
\texttt{23Q-p-Y8} & 184 & 2{,}353{,}250 & 1.2.1 & 247.8 & 394 & 12 & 4 & 1{,}204 & 51.4 \\
\texttt{23R-p-Y8} & 184 & 2{,}364{,}374 & 1.2.1 & 158.6 & 293 & 10 & 4 & 1{,}174 & 46.9 \\
\texttt{23S-p-Y8} & 184 & 2{,}353{,}250 & 1.2.1 & 431.6 & 643 & 6 & 2 & 1{,}114 & 47.1 \\
\texttt{23T-p-Y8} & 184 & 2{,}364{,}374 & 1.2.1 & 164.9 & 312 & 12 & 4 & 1{,}204 & 48.4 \\
\texttt{23U-p-Y8} & 184 & 2{,}353{,}250 & 1.2.1 & 242.6 & 401 & 12 & 4 & 1{,}204 & 53.5 \\
\texttt{23V-p-Y8} & 184 & 2{,}353{,}250 & 1.2.1 & 327.4 & 483 & 13 & 4 & 1{,}219 & 52.1 \\
\texttt{23W-p-Y8} & 184 & 2{,}196{,}265 & 1.2.1 & 68.1 & 139 & 0 & 0 & 1{,}024 & 9.4 \\
\texttt{23X-p-Y8} & 184 & 2{,}196{,}265 & 1.2.1 & 72.2 & 144 & 0 & 0 & 1{,}024 & 9.6 \\
\texttt{23Y-p-Y8} & 184 & 2{,}219{,}145 & 1.2.1 & 14.5 & 48 & 0 & 0 & 1{,}024 & 4.7 \\
\midrule
Total & & & & 4{,}800.7 & 10{,}336 & 218 & & 31{,}942 & 1{,}149.9 \\
\bottomrule
\end{tabular}
\end{table}

\subsection{Resources}\label{sec:resources}
All cloud computations ran on Amazon Web Services Graviton4 spot instances in
the region eu-north-1, each with $96$ vCPUs: one c8g.24xlarge for the
enumeration and three m8g.24xlarge (one for the new instances, two in
succession for NMH's instances). The remainder of the $n=21$ enumeration,
$5{,}259$ jobs, ran on the machine later used for NMH's instances, after the
enumeration machine was lost to a spot interruption. Solver time was about
$1{,}660$ hours for the catalogues, $2{,}604$ hours for the new instances and
$1{,}150$ hours for NMH's instances, about $5{,}400$ hours in all. These are
sums of per-job or per-cube solver wall-clock times. At times the number of
workers exceeded the number of vCPUs, so they are not CPU-hours. The cost,
metered from public spot prices, was about US\$39.
% TODO: replace by the final figure from the AWS bill. The meter's last reading
% (results/aws_cost.json, 10:02 UTC on 5 Nov) was US$38.65 with the last machine still
% running until its scheduled power-off.
For comparison, NMH report about $8$ days of CaDiCaL time on average for their
$25$ instances with $Y_8$, on one core each~\cite[\S5.2.3]{NMH22}; their logs give
$4{,}801$ hours for the $27$ logged instances, without proofs.

% =====================================================================
\section{Trust base}\label{sec:trust}

The proof checker covers the unsatisfiability of every formula that the proof
uses. The following parts are \emph{not} checked by it. A reader who accepts
these items and the checked computations can accept
Theorem~\ref{thm:main}.

\begin{enumerate}[label=\textbf{T\arabic*.},leftmargin=*]
\item \textbf{The mathematical argument.} Lemmas~\ref{lem:pivot}--\ref{lem:h6g7},
  the reduction of Section~\ref{sec:reduction}, and
  Propositions~\ref{prop:pivotdec} and~\ref{prop:complete} and
  Lemmas~\ref{lem:sound} and~\ref{lem:cover}. They are short and checked by
  hand only; nothing is formalised.

\item \textbf{One theorem from the literature.} The uniqueness of the
  $\TT{6}$-free tournament on $26$ vertices (Fact~\ref{fact:st26}). Everything
  else we take from the literature is re-derived: $R(5) \le 14$ by our
  $\TT{5}$-free lists, $R(6) \le 28$ from it by Lemma~\ref{lem:double},
  NMH's classifications on $23$--$25$ vertices by our certified enumeration, and
  NMH's $28$ SAT results by our certified runs. Lemma~\ref{lem:h6g7} has a
  short proof. Fact~\ref{fact:st26} could be re-derived with the same
  pipeline: for $n=26$ the degree bounds force $d=13$, and there is one
  $\TT{5}$-free tournament on $12$ vertices and one on $13$, so the enumeration
  is a single pivot job. We have not run it.
  % TODO (author): decide whether to run and certify the single n = 26 job before submission.

\item \textbf{The $\TT{5}$-free lists.} Their completeness (Section~\ref{sec:tt5})
  rests on our extension program and on \texttt{labelg}; it is not certified.
  It was cross-checked against \texttt{gentourng} for $n\le 10$ and against
  NMH's files for $n \le 13$.

\item \textbf{Soundness of the enumeration formulas.} That every
  $\TT{6}$-free completion of a job with out-degrees at most $d$ satisfies the
  job formula. This is Lemma~\ref{lem:sound} plus the correctness of the
  sequential counter, \emph{as implemented} in \path{src/ttenc.h} and
  \path{src/pivenum.cpp}. These formulas were not checked clause by clause.
  The evidence is the brute-force tests, a test that the enumerator's
  $\TT{k}$ clauses equal those of \path{src/encode.c}, the reproduction of the
  $\TT{5}$-free lists for $n=10,\dots,13$ by pivot enumeration, and the exact
  agreement of the final lists with two independent sources.

\item \textbf{Enumeration bookkeeping.} That the job list is complete, that
  each blocking clause is built from the recorded model, and that each recorded
  model agrees with the fixed arcs of its job. The certification step checks
  that recorded models are $\TT{6}$-free but not that they agree with the fixed
  arcs; the enumerator builds them from those arcs.
  % NOTE (author): a one-line check in scripts/certify_enum.py would remove this item.

\item \textbf{Catalogue analysis.} The $Y_8$-containment search and the
  identification of $\Ostar$ (Python, \path{src/tourn.py}); and the property
  of \texttt{labelg} that its output is a relabelling of its input.

\item \textbf{Soundness of the five new formulas.} They were produced by our
  encoder in its $3$-cycle mode ($L=1$). The clause checker of
  Section~\ref{sec:stage4enc} handles only direct encodings and was \emph{not}
  run on them. Their soundness therefore rests on Lemma~\ref{lem:sound} and on
  \path{src/encode.c} implementing it, which is supported by the brute-force
  tests (which include $L=1$). Also trusted: the construction of the instance
  matrices (\path{src/instances.py}) from \path{data/I24_X.txt} and
  \path{data/O21star.txt}, and the identification of these files with
  $\cl{24X}$ and $\Ostar$ by canonical form.
  % TODO (author): extending src/cnf_sound.c to 3-cycle variables and running it on the
  % five CNFs (a few seconds each) would move most of this item into the checked part.

\item \textbf{Soundness of NMH's formulas.} This rests on the clause checker
  \path{src/cnf_sound.c} and on the structure check of the inputs, not on NMH's
  encoder.

\item \textbf{The LRAT checker.} \texttt{lrat-check} is an unverified C
  program. A formally verified checker was not used.

\item \textbf{The streaming harness and cube bookkeeping.} That
  \path{scripts/stream_solve.py} gives the checker the same formula $F \wedge c$
  as the solver and reads its verdict correctly; that
  \path{scripts/run_cubes.py} records the literals of the cube that was
  actually solved; and the structural split check and leaf tally in
  \path{scripts/cnc_summary.py}. (That script compares every child with its
  parent's literals, but it does not compare a level-$0$ record with the cube
  file; the two agree by construction in \path{run_cubes.py}.) The level-$0$ cover itself is checked by
  LRAT.

\item \textbf{Proofs are not kept.} Each LRAT proof was checked once, while it
  was produced, and then discarded. The published record consists of per-cube
  and per-job verdicts. An independent re-check means re-running the solver,
  about $3{,}750$ solver-hours for the $33$ pivot instances and about
  $1{,}660$ for the catalogues.

\item \textbf{Faults.} The runs used ordinary cloud hardware. A fault in a
  solver run is caught by the checker. A fault in the checker or in the
  bookkeeping is not, except where the recomputation of the summaries on a
  second machine would expose it.
\end{enumerate}

The following are \emph{not} trusted: CaDiCaL (its unsatisfiability claims are
checked and its models are re-checked), NMH's encoder, NMH's runs and
classifications, and the catalogue files of~\cite{LLHrepo}, which serve only as
a cross-check.

% =====================================================================
\section{Data and code availability}\label{sec:data}

The code, the input data, the result files and the notes are in a public
repository (TODO: URL).
% TODO: repository URL and an archived snapshot with DOI (e.g. Zenodo) will be added.
An archived snapshot will be deposited with a DOI (TODO). The repository
contains, for every enumeration job, the completion record and the certificate
record; for every cube, its literals, status, checker verdict and times; the
per-instance summaries from the cloud machines and from the second machine;
and the SHA-256 manifest of NMH's instances. It does not contain LRAT proofs
(Section~\ref{sec:trust}, item T11) or third-party files. NMH's repository and
the catalogue files of~\cite{LLHrepo} are fetched by
\path{scripts/fetch_refs.sh}; the SHA-256 values of the files we used are in
\path{refs/SHA256SUMS.txt}. The SHA-256 values of the input matrices of the
five new instances are in Appendix~\ref{app:data}. The main programs are:
{\raggedright
\begin{itemize}
\item enumeration: \path{src/pivenum.cpp}, \path{src/ttenc.h},
  \path{scripts/pivot_enum.py}, \path{scripts/certify_enum.py},
  \path{scripts/certify_first.py}, \path{scripts/stage12_check.py};
\item encoding and checking of formulas: \path{src/encode.c},
  \path{src/instances.py}, \path{src/cnf_sound.c},
  \path{scripts/stage4_instances.py};
\item cube-and-conquer: \path{scripts/make_cubes.py},
  \path{scripts/run_cubes.py}, \path{scripts/stream_solve.py},
  \path{scripts/cnc_summary.py}, and the cloud drivers in
  \path{scripts/cloud/};
\item tournament utilities: \path{src/tourn.py}.
\end{itemize}}

% =====================================================================
% DRAFT — the author must review and rewrite this section
\section{Use of AI tools}\label{sec:ai}

The code in the repository, the preparation and monitoring of the cloud runs,
the working notes and the initial draft of this text were produced with an AI
assistant (Claude, Anthropic) working under the author's direction. The
commit messages of the repository name the assistant. The author chose the problem
and approved each computational stage and its budget. The AI system is not an
author of this paper. The author takes responsibility for its content.

% =====================================================================
\section*{Acknowledgements}
We thank David Neiman, John Mackey and Marijn Heule for publishing their code,
inputs and solver logs, and Evan Lohn, Chris Lambert and Marijn Heule for
publishing their catalogue files. This work would not have been possible
without them. We also thank the developers of CaDiCaL, drat-trim and nauty for
making their tools freely available.

\bibliographystyle{amsplain}
\bibliography{refs}

% =====================================================================
\appendix
\section{Data}\label{app:data}

Adjacency matrices have entry $1$ in row $i$, column $j$ if and only if
$i \to j$ (vertices numbered from $0$). These are the matrices of
\path{data/Y8.txt}, \path{data/H6.txt} and \path{data/W7.txt}; $Y_8$ is the
matrix of \cite[Fig.~5]{NMH22}.
{\small\setlength{\arraycolsep}{2pt}\renewcommand{\arraystretch}{0.9}
\[
Y_8:\ \begin{matrix}
0&0&1&1&1&1&1&1\\
1&0&0&0&1&0&1&1\\
0&1&0&0&1&1&0&1\\
0&1&1&0&0&1&1&0\\
0&0&0&1&0&1&0&1\\
0&1&0&0&0&0&1&1\\
0&0&1&0&1&0&0&0\\
0&0&0&1&0&0&1&0
\end{matrix}
\qquad
H_6:\ \begin{matrix}
0&1&1&1&1&1\\
0&0&1&1&1&1\\
0&0&0&1&1&1\\
0&0&0&0&1&0\\
0&0&0&0&0&1\\
0&0&0&1&0&0
\end{matrix}
\qquad
G_7:\ \begin{matrix}
0&1&1&1&1&1&1\\
0&0&1&1&1&1&1\\
0&0&0&1&1&0&0\\
0&0&0&0&1&1&0\\
0&0&0&0&0&1&1\\
0&0&1&0&0&0&1\\
0&0&1&1&0&0&0
\end{matrix}
\]}

The tournament $\Ostar$, as labelled in \path{data/O21star.txt}, in nauty's
digraph6 format (the adjacency matrix row by row, six bits per character):
\begin{verbatim}
&T^}?@Z\GEmsQMySCw[xQdZPLFhCtfDQViPXZB_N|ZQSenHEMsakhTUgkUyOxdqIdxCw]dKX{@}?
\end{verbatim}

SHA-256 of the input files of the five new instances:
\begin{center}\footnotesize
\begin{tabular}{@{}ll@{}}
\path{data/I24_A.txt} & \texttt{7e027dfd8cde62f9b866fb05f8e8225449c3e218dd23efa361fe8e59b3b518d1} \\
\path{data/I24_B.txt} & \texttt{4bbbdb4138e61a164dd583d4a1c4a750504b54a15fe95a230aac5644c7ba37b3} \\
\path{data/I24_C.txt} & \texttt{eb94d398672fd48006d1319c01c48c0cdafcb23ebb22bdebe3e111368b4dc87c} \\
\path{data/I24_D.txt} & \texttt{fb94bdde12356261fc2a698c4fe37b8b64bacdecdcae81f68310a6aa6b7a8f2c} \\
\path{data/I24_E.txt} & \texttt{73a635a9bd95129cc13d50f4b34e311ad45abf97a405d81995653f045cfe7c35} \\
\path{data/O21star.txt} & \texttt{1823ac5f5d77c1244622be79d66743100a8306c4fc554b6f3db9f09d1a6f641b} \\
\end{tabular}
\end{center}
% TODO: the SHA-256 values of the five Stage 3 CNFs were not recorded on the server;
% the clause and literal counts in results/stage3-ABCDE.log are the only fingerprint.

\end{document}
